\section{问题分析}

\subsection{总体分析}

四个问题的物理系统完全相同：一张光储微网与一个外部电网相连，储能设备的容量、功率、效率与运行区间由题面给定，光伏先供负载、盈余充入储能、不足由外网补足。因此四问的差异不在物理规律，而在\textbf{决策者知道什么、什么时候决定、按什么规则结算}三件事上。据此本文把四个问题统一抽象为四个维度上的参数化：

\begin{enumerate}[leftmargin=2em,itemsep=1pt]
  \item \textbf{时间维度}（决策时刻集合）：问题一、二每天只有一个决策时刻；问题三与链 \emph{4-3} 有 0:00、6:00、12:00、18:00 四个决策时刻，后续时刻可修改先前的决定。
  \item \textbf{信息维度}（各层信息集）：问题一、二的计划层被视为完全信息；问题三的计划层只知 0:00 发布的整点预报，调整层另可见当天已实现时段；链 \emph{4-2}、\emph{4-3} 在 0:00 时对当日实时电价一无所知，只能用历史价格。
  \item \textbf{结算维度}（超额与不足如何计价）：问题一、二只有“按计划量计费”与“紧急购电按 5 倍价”两条规则；问题三与链 \emph{4-3} 另有 $0.5\times/1.5\times$ 双向分段结算。
  \item \textbf{价格情景维度}：电价是逐日重复的单日曲线（附件 1），还是逐日波动且当日不可预知的实际序列（附件 4）。
\end{enumerate}

这一抽象带来两个好处：四个问题共用一套平衡方程、储能动态、容量与功率约束和目标函数，模型规模与正确性可用同一套方法核验；同时“难点”被定位到具体维度上，“是否需要更多时刻的预报”“要不要做随机优化”“日边界为什么总贴储电量下限”这类问题都有了明确的讨论对象。技术路线如图~\ref{fig:roadmap}。

\begin{figure}[htbp]
\centering
\begin{tikzpicture}[font=\footnotesize,
  box/.style={draw,rounded corners=2pt,align=center,inner sep=3pt,minimum height=6.5mm},
  core/.style={box,fill=black!7,thick},
  dim/.style={box,fill=black!2,dashed},
  ar/.style={->,>=stealth,thick}]
\node[dim] (d1) at (0,2.0) {时间维度\\决策时刻集合};
\node[dim] (d2) at (2.9,2.0) {信息维度\\各层信息集};
\node[dim] (d3) at (5.8,2.0) {结算维度\\超额/不足计价};
\node[dim] (d4) at (8.7,2.0) {价格情景\\附件 1 / 附件 4};
\node[core] (lp) at (4.35,0.6) {统一逐时段确定性线性规划内核\\平衡方程 $\cdot$ 储能动态 $\cdot$ 容量与功率约束 $\cdot$ 费用目标};
\node[box] (p1) at (0,3.4) {问题一\\单日计划};
\node[box] (p2) at (2.9,3.4) {问题二\\全年计划 + 紧急购电};
\node[box] (p3) at (5.8,3.4) {问题三\\计划 + 调整 + 分段结算};
\node[box] (p4) at (8.7,3.4) {问题四\\两条链并行};
\draw[ar] (d1) -- (lp); \draw[ar] (d2) -- (lp); \draw[ar] (d3) -- (lp); \draw[ar] (d4) -- (lp);
\draw[ar] (lp) -- (p1); \draw[ar] (lp) -- (p2); \draw[ar] (lp) -- (p3); \draw[ar] (lp) -- (p4);
\end{tikzpicture}
\caption{技术路线：在四个维度上参数化的统一线性规划内核}
\label{fig:roadmap}
\end{figure}

\subsection{各问题的关键难点}

\textbf{问题一}的输入（附件 1 的电价、负载与光伏预测）均视为确定值，约束中真正起作用的是储能：容量上下限、两侧功率上限、往返效率 81\%，以及“0:00 与 24:00 储电量相等”的日周期条件。经济上这是典型的跨时段套利：储能把电能从低价时段搬到高价时段，搬运损耗表现为往返效率小于 1。需要论证的是两件事：是否需要二元变量强制“不同时充放电”；光伏盈余超过储能吸收能力时功率平衡如何闭合。

\textbf{问题二}需处理三点新结构。其一，日周期约束被\textbf{跨日状态连续性}取代——题面只在问题一要求端点储量相同，把该约束逐日施加会人为割裂跨日搬运电能的可能性，本文用逐日解耦的对照量化这一差别。其二，\textbf{终端条件必须显式声明}：题面未规定决策期末储电量，若允许自由终止，模型会把期末库存“变卖”以压低账面费用，本文保留终端自由但把该效应单独计价。其三，\textbf{紧急购电在完全信息下不会被激活}：既然 0:00 已知当天实际负载与光伏、且购电无功率上限，任何“先欠着、事后按 5 倍价补”都可事先改为按正常价买入同样电量而严格更省。该结论是模型内的推论而非实测断言，也不意味着机制可删——一旦引入购电上限，紧急购电立即被激活。

\textbf{问题三}的关键变化是\textbf{信息结构}：附件 3 只给整点预报，而决策与结算以 10 分钟为粒度。由此产生两个必须处理的环节：\emph{降尺度}（本文取整点锚定线性插值，该口径不保证块内能量守恒，其影响可能大于调整机制本身，故作为独立消融项量化）与\emph{非预期性}（0:00 计划只能用 0:00 预报，调整只能用该时刻已发布预报与已实现状态，不得使用未来实际值）。因此模型被组织为三层顺序向前递推，而非能“看到未来”的联合优化。此外，由于决策层用预报、结算层用实际光伏，预报误差必然造成缺口；调整降低了紧急购电却引入偏差违约金，其净效果只能用同一口径下的费用分解回答。

\textbf{问题四}的实质是\textbf{价格可观测性}的改变。“实时”意味着 0:00 制定当天计划时当天 144 个电价点尚未发生；若直接把附件 4 当作已知输入，就等于假设电价可完全预知。本文明确采用“决策层用历史价格构造的预测价、结算层一律用实际价”的口径，并把二者差额定义为“价格预报误差的代价”单独报告。由此两条链的性质不同：链 \emph{4-2} 中价格不确定性只改变费用、不产生能量缺口；链 \emph{4-3} 中缺口来自光伏预报误差，价格不确定性则同时影响计划层、调整层的决策与最终账单。两链费用差额可完全分解为计划口径差、偏差结算与紧急购电三项。同时，两链都采用“当天决策、终端自由”的组织方式，日边界储电量会贴住下限——这究竟是普遍性质还是特定结构的产物，需要改变前瞻深度的对照实验来回答。
